\section{Conclusion}
\label{sec:conclusion}

Joggle treats compiler extension as typed computation over a shared graph.
Compiler functions provide one surface for semantics, analysis,
transformation, conversion, and artifact generation. Mods make capabilities
explicit units of ownership and composition alongside hierarchical program IR.
Revisions, observed dependencies, effect propagation, transactions, and cached
execution plans make repeated compilation change-proportional. Together, these
mechanisms connect three properties that compiler infrastructures usually
expose separately: a predictable way to write an extension, a controlled
boundary in which it evolves, and selective work after it changes. The
evaluation tests these properties through executable extension completion,
paired patch footprint, and reactive update cost while measuring generated
artifacts independently.
